\documentclass[11pt,a4paper,sans,english]{moderncv}        % possible options include font size ('10pt', '11pt' and '12pt'), paper size ('a4paper', 'letterpaper', 'a5paper', 'legalpaper', 'executivepaper' and 'landscape') and font family ('sans' and 'roman')
\moderncvstyle{casual}                             % style options are 'casual' (default), 'classic', 'oldstyle' and 'banking'
\moderncvcolor{blue}                               % color options 'blue' (default), 'orange', 'green', 'red', 'purple', 'grey' and 'black'
%\nopagenumbers{}                                  % uncomment to suppress automatic page numbering for CVs longer than one page
\usepackage[utf8]{inputenc}                       % if you are not using xelatex ou lualatex, replace by the encoding you are using
\usepackage[scale=0.80,a4paper]{geometry}
\usepackage{babel}
\usepackage{mathtools}
\usepackage{amssymb}
\usepackage{needspace}
%----------------------------------------------------------------------------------
%            personal data
%----------------------------------------------------------------------------------
\firstname{}
\familyname{Davide Tornotti}
%\title{Resumé title}                               % optional, remove/comment the line if not wanted
\address{University of Milano-Bicocca}{Department of Physics}{Italy}         % optional, remove/comment the line if not wanted; the "country" arguments can be omitted or provided empty
%\mobile{mobile number}                          % optional, remove/comment the line if not wanted
%\phone{phone number}                           % optional, remove/comment the line if not wanted
%\fax{fax number}                             % optional, remove/comment the line if not wanted
\email{d.tornotti@campus.unimib.it}                               % optional, remove/comment the line if not wanted
%\homepage{home page}                         % optional, remove/comment the line if not wanted
%\extrainfo{additional information}                 % optional, remove/comment the line if not wanted
%\photo[68pt][0.4pt]{/Users/davidetornotti/Desktop/DAVIDE/Foto_profilo.jpg}
%\quote{some quote}                                 % optional, remove/comment the line if not wanted
%

\newcommand{\pubentry}[4]{%
	#1 #2, #3, #4\\[1mm]
}


%----------------------------------------------------------------------------------
%   ORCID: replace the placeholder below with your ORCID iD (format 0000-0000-0000-0000)
%----------------------------------------------------------------------------------
\newcommand{\myorcid}{0009-0001-3388-8742}

\begin{document}
%-----       resume       ---------------------------------------------------------
\makecvtitle

\section{Personal Information}
\cvitem{Email:}{d.tornotti@campus.unimib.it}
\cvitem{Link:}{\url{https://en.unimib.it/davide-tornotti}}
\cvitem{ORCID:}{\href{https://orcid.org/\myorcid}{\myorcid}}
\cvitem{Nationality:}{Italian}

\section{Scientific Interests}
\cvitem{}{Cosmic web, circumgalactic and intergalactic medium, Ly$\alpha$ emission and radiative transfer, galaxy evolution in different environments, cosmological simulations, large-scale structure, integral-field spectroscopy}

\section{Education}
\cventry{2023--}{PhD in Astrophysics (expected September 2027)}{Università degli Studi di Milano-Bicocca}{Italy}{}{Title: Unveiling the cosmic structures with ultra-deep spectroscopic data} 
\cventry{2021--2023}{MSc in Physics}{Università degli Studi di Milano}{Italy}{110/110 cum laude}{Thesis: Reliability of the Ly$\alpha$ line as a virial broadening estimator for quasars at high redshift} 
\cventry{2018--2021}{BSc in Physics}{Università degli Studi di Milano}{Italy}{110/110 cum laude}{Thesis: Determination of the supermassive black holes mass in a sample of Type 1 AGN from the VIPERS survey}  

\section{Awards and Fellowships}
\cventry{2026}{Enrico Fermi Fellowship}{John Templeton Foundation}{}{}{}

\section{Publications}
\cvitem{}{7 refereed papers (4 first-author, including \textit{Nature Astronomy} and \textit{ApJL}); 1 public software package; 1 white paper. Full list: \href{https://ui.adsabs.harvard.edu/search/q=orcid\%3A\myorcid}{NASA/ADS}.}

\subsection{Refereed -- first author}
\cvitem{1.}{\textbf{Tornotti, D.}, Fumagalli, M., Fossati, M., et al. 2026, A\&A, 715, A59. \textit{Cosmic web Ly$\alpha$ emission in a sample of overdense regions.}}
\cvitem{2.}{\textbf{Tornotti, D.}, Fossati, M., Fumagalli, M., et al. 2025, A\&A, 704, A201. \textit{Bayesian luminosity function estimation in multidepth datasets with selection effects: a case study for $3<z<5$ Ly$\alpha$ emitters.}}
\cvitem{3.}{\textbf{Tornotti, D.}, Fumagalli, M., Fossati, M., et al. 2025, ApJL, 980, L43. \textit{The MUSE Ultra Deep Field: a 5\,Mpc stretch of the $z\approx4$ cosmic web revealed in emission.}}
\cvitem{4.}{\textbf{Tornotti, D.}, Fumagalli, M., Fossati, M., et al. 2025, Nature Astronomy, 9, 577--588. \textit{High-definition imaging of a filamentary connection between a close quasar pair at $z=3$.}}

\needspace{5\baselineskip}
\subsection{Refereed -- co-author}
\cvitem{5.}{Mukherjee, T., Li, Z., \ldots, \textbf{Tornotti, D.}, et al. 2026, A\&A, 710, A331. \textit{BLUEPRINT: Blue-dominant Lyman-$\alpha$ emission as evidence of gas inflow in ultra-low-mass galaxies at $z=3$.}}
\cvitem{6.}{Galbiati, M., \textbf{Tornotti, D.}, Fumagalli, M., et al. 2026, A\&A, 709, A233. \textit{MUSE Analysis of Gas around Galaxies (MAGG) VII. Emission-line galaxies near strong blended Ly$\alpha$ absorption systems at $z\gtrsim3$.}}
\cvitem{7.}{Pensabene, A., Galbiati, M., Fumagalli, M., \ldots, \textbf{Tornotti, D.}, et al. 2025, A\&A, 696, A33. \textit{The MUSE Ultra Deep Field (MUDF) VII. Probing high-redshift gas structures in the surroundings of ALMA-identified massive dusty galaxies.}}

\subsection{Software}
\cvitem{8.}{Fossati, M., \& \textbf{Tornotti, D.} 2025, Zenodo, v1.1.0, \href{https://doi.org/10.5281/zenodo.14710518}{doi:10.5281/zenodo.14710518}. \textit{SHINE (Spectral Highlighting and Identification of Emission)}: Python package to detect and visualise emission features in integral-field spectroscopic data (\href{https://github.com/matteofox/SHINE}{GitHub}).}


\subsection{Other publications (non-refereed)}
\cvitem{9.}{Goovaerts, I., Beckett, A., Hayes, M., \ldots, \textbf{Tornotti, D.}, et al. 2026, arXiv:2606.04136. \textit{HST's Deep Blue: extremely deep UV imaging to reveal the contributors to reionization} (white paper).}

\section{Leadership and Collaborations}
\cvitem{2025--}{\textbf{Wide-field Spectroscopic Telescope (WST)}: lead contributor to the Cosmic Web science case (Cosmic Web Tiger Team)}
\cvitem{2024--}{\textbf{MUSE Ultra Deep Field (MUDF) survey}: lead of the MUSE analysis for intergalactic medium studies}
\cvitem{2025--}{\textbf{MUSERePipe}: co-lead of the archival project to build a homogeneous, science-ready collection of the deepest MUSE observations}
\cvitem{2025--}{\textbf{MIRACLES survey}: team member contributing to the HSC analysis for intergalactic medium studies}

\section{Observing Programmes (PI and Co-I)}
\cvitem{2026}{\textbf{PI}, ESO VST/OmegaCAM, 47\,h (117.2B2G.001): \textit{Tracing the large-scale environment of the MUDF filaments with LBGs}}
\cvitem{2026}{\textbf{Co-I}, ESO VLT/MUSE, 40\,h (117.29XU): \textit{Detailed characterization of the cosmic web filaments at $z=3$}}

\section{Mentoring and Teaching}
\cvitem{2024--}{Co-supervisor of 5 Bachelor and Master theses, University of Milano-Bicocca}
\cvitem{2025}{Invited hands-on session on volumetric data analysis in astrophysics, Imaging25, Varenna, Italy}
\cvitem{2025--}{Teaching Assistant, Laboratory of data acquisition, MSc in Astrophysics and Space Physics, University of Milano-Bicocca}
\cvitem{2025--}{Tutor, Introduction to galaxies, MSc in Astrophysics and Space Physics, University of Milano-Bicocca}

\needspace{6\baselineskip}
\section{Service and Public Engagement}
\cvitem{2025--}{Referee for \textit{Astronomy \& Astrophysics}}
\cvitem{Jun., 2024}{Member of the LOC, \textit{What Matter(s) Around Galaxies 2024}, Varenna, Italy}
\cvitem{2024--}{Outreach activities and planetarium sessions at University of Milano-Bicocca}

\section{Research Visits}
\cventry{30-31/10/25}{Visitor}{Kavli Institute for Cosmology}{University of Cambridge}{Cambridge}{}
\cventry{28/10/25}{Visitor}{The Royal Observatory}{University of Edinburgh}{Edinburgh}{}{}
\cventry{27/10/25}{Visitor}{University of Durham}{Durham}{}{}
\cventry{05--08/25}{Visitor}{Max Planck Institute for Astrophysics}{Germany}{}{}

\section{Talks, Seminars, Conferences}
\cvitem{}{24 talks and seminars since 2024, of which 8 invited.}
\cventry{Oct., 2026}{Filaments in the multi-scale Universe}{Kavli IPMU}{Japan}{}{Unveiling the cosmic web through Ly$\alpha$ emission}
\cventry{Jul., 2026}{ESO Expanding Horizon}{ESO}{Germany}{The $z\sim2-5$ cosmic web ecosystem: a revolutionary vantage point with WST}{}
\cventry{Jun., 2026}{Bicoq Conference: from gravity to particles}{Milan}{Italy}{Mapping the Cosmic Web in Ly$\alpha$ emission}{}
\cventry{Jun., 2026}{RASCAS User Meeting}{Lyon}{France}{Modeling the diffuse Ly$\alpha$ emission of the Cosmic Web:
	RASCAS post-processing of COLIBRE simulations}{}
\cventry{Nov., 2025}{The multiscale environment of AGN across cosmic time}{Ringberg}{Germany}{Invited}{Mapping the Cosmic Web in Ly$\alpha$ emission Around Quasars and Non-AGN Galaxies}
\cventry{Oct., 2025}{Galaxies Discussion Group}{Kavli Institute for Cosmology}{Cambridge}{UK}{Mapping the Cosmic Web in Ly$\alpha$ emission Around Quasars and Non-AGN Galaxies}
\cventry{Oct., 2025}{Coffee Talk}{The Royal Observatory}{Edinburgh}{UK}{Mapping the Cosmic Web in Ly$\alpha$ emission Around Quasars and Non-AGN Galaxies}
\cventry{Oct., 2025}{Seminar}{University of Durham}{Durham}{UK}{Mapping the Cosmic Web in Ly$\alpha$ emission Around Quasars and Non-AGN Galaxies}
\cventry{Oct., 2025}{IASF Colloquium Series}{IASF-INAF}{Milan}{Invited}{Mapping the Cosmic Web in Ly$\alpha$ emission Around Quasars and Non-AGN Galaxies}
\cventry{Sept., 2025}{Imaging25}{Varenna}{Italy}{}{
	Invited talk: Tomography imaging and data analysis techniques for volumetric data: applications in astrophysics \\
	Invited hands-on session: Methods and applications of volumetric data analysis in astrophysics}
\cventry{Sept., 2025}{WEAVE-QSO Meeting}{Milano}{Italy}{Invited}{From emission to absorption: probing circumgalactic and intergalactic gas}
\cventry{Aug., 2025}{Feedback with Friends}{Ascona}{Switzerland}{Invited}{Mapping the Cosmic Web in Ly$\alpha$ emission Around Quasars and Non-AGN Galaxies}
\cventry{Jul., 2025}{CGM Meeting}{ESO}{Garching}{}{Mapping the cosmic web in Ly$\alpha$ emission}
\cventry{Jun., 2025}{OPINAS Seminar}{MPE}{Garching}{Invited}{Mapping the cosmic web in Ly$\alpha$ emission}
\cventry{Jun., 2025}{EAS}{SS7a: WST: the Wide-field spectroscopic facility}{Cork}{Ireland}{The cosmic web in Ly$\alpha$ emission: from MUSE to WST.}
\cventry{Jun., 2025}{Cosmology Seminar}{MPA}{Garching}{}{Unveiling the cosmic web through Ly$\alpha$ emission:  an ultra-deep observation with MUSE}
\cventry{Apr., 2025}{Café-Club}{LAM}{Online}{Invited}{Unveiling the cosmic web through Ly$\alpha$ emission:  an ultra-deep observation with MUSE}
\cventry{Apr., 2025}{Escape of Lyman radiation from galactic labyrinths}{OAC}{Crete}{}{Unveiling the cosmic web through Ly$\alpha$ emission:  an ultra-deep observation with MUSE}
\cventry{Mar., 2025}{WST: Surveying the Universe in the 2040's and beyond}{INAF}{Naples}{}{Unveiling the cosmic web through Ly$\alpha$ emission: from MUSE to WST}
\cventry{Jan., 2025}{Cosmolunch Colloquium}{HUJI}{Online}{Invited}{Unveiling the cosmic web through Ly$\alpha$ emission: an ultra-deep observation with MUSE}
\cventry{Nov., 2024}{A decade of discoveries with MUSE and beyond}{ESO}{Germany}{}{High-definition imaging of an extended filament linking two active quasars at cosmic noon}
\cventry{Oct., 2024}{Beyond the Edge of the Universe}{Sintra}{Portugal}{}{Unveiling the cosmic web through Ly$\alpha$ emission: an ultra-deep observation with MUSE}
\cventry{Jun., 2024}{What Matter(s) Around Galaxies 2024}{Varenna}{Italy}{}{A high-resolution view of the cosmic web in emission at high redshift}
\cventry{Jun., 2024}{Where the Circum Galactic Medium meets the galaxy environment}{IFPU}{Italy}{}{A high-resolution view of the cosmic web in emission at high redshift}

\section{Schools}
\cvitem{Oct., 2024}{Artificial Intelligence and Modern Physics: a two-way Connection, Monopoli, Italy}
\cvitem{Jan., 2024}{Machine and Deep Learning, INAF-IASF, Italy}

\section{Languages}
\cvitemwithcomment{Italian}{Mother Tongue}{}
\cvitemwithcomment{English}{Fluent}{}


\needspace{4\baselineskip}
\section{Volunteering}
\cvitem{2024--}{Climbing Instructor, CNSASA, Club Alpino Italiano, Italy}
\cvitem{2018--}{Responsible for Makalu Rock Cube Climbing Gym, Club Alpino Italiano, Cameri, Italy}
\cvitem{2018--}{Board Member, CAI Cameri Section, Italy}


\end{document}
